% Figure F4 --- the hallucination adverse-event lifecycle (Section 8).
% \resizebox to \textwidth in main.tex.
\begin{tikzpicture}[
  st/.style={draw, rounded corners, align=center, minimum height=1.15cm,
             text width=2.05cm, font=\scriptsize, fill=black!3},
  arr/.style={-{Latex[length=2mm]}, thick, black!70},
  fb/.style={-{Latex[length=1.8mm]}, dashed, black!55},
  lbl/.style={font=\scriptsize\itshape, text=black!55},
  node distance=6mm
]
\node[st] (s1) {\textbf{Instance}\\a single failure};
\node[st, right=of s1] (s2) {\textbf{Incidence}\\recurrence tracked over time};
\node[st, right=of s2] (s3) {\textbf{Causality \& cascade}\\traced to a cause; chains mapped};
\node[st, right=of s3] (s4) {\textbf{Mitigation}\\intervention applied};
\node[st, right=of s4] (s5) {\textbf{Efficacy}\\rate-change across a dated boundary};
\node[st, right=of s5] (s6) {\textbf{Prevention}\\moved upstream};
\draw[arr] (s1) -- (s2);
\draw[arr] (s2) -- (s3);
\draw[arr] (s3) -- (s4);
\draw[arr] (s4) -- (s5);
\draw[arr] (s5) -- (s6);
\draw[fb] (s5.south) to[out=-90,in=-90,looseness=0.35]
     node[below,lbl]{reduces future incidence} (s2.south);
\node[lbl] at ($(s1.north)!0.5!(s6.north)+(0,6mm)$)
     {each step is a query over the linked record (Section~\ref{sec:datamodel})};
\end{tikzpicture}
